\section{Preliminaries}\label{pre}
% ------------------------------------------------------------------------------
\paragraph{Assumptions}
For simplicity, we first assume that periods of sex work
can be precisely and instantaneously delineated.
That is, we assume clear boundaries between formal sex work and transactional sex,
and that surveys predominantly reach those actively engaged in sex work,
despite some surveys including anyone who
bought/sold sex anytime in the past year, etc. \cite{Scorgie2012}.
We further assume that survey participants
can accurately report how long they have bought/sold sex
--- \ie perfect recall, no social desirability bias.
Finally, we assume that all populations are at equilibrium,
\eg rates of entering/exiting sex work are stable over time.
% ------------------------------------------------------------------------------
\paragraph{Who exactly are we talking about?}
Before continuing, it is helpful to distinguish two populations:
individuals who will ever sell/buy sex in their lifetime, and
individuals who are actively selling/buying sex.
The former "source population" is the sampling frame for method~1,
and tends to include more individuals who sell/buy sex for shorter durations;
the latter "active population" is the sampling frame for methods 2 and 3,
and tends to include more individuals who sell/buy sex for longer durations
(Figure~\ref{fig:diag.pop}).
Since "sex work" strata in compartmental models reflect the active population,
we might assume that rates of exit from sex work
should be derived from the active population.
% SM: should the word be "could" instead of "should"?
% JK: I think I lean towards "should" because I am imagining
%     someone trying to determine the "right" way to do it,
%     whereas "could" suggests there different options e.g. with pros/cons
However, this is not the case, because
these exit rates apply to all individuals \emph{entering} sex work,
which reflects the source population.%
\footnote{An analogous example: % TODO: consider moving to main text
  the mean duration of an infection should be characterized by
  a cohort of individuals with \emph{incident} infections,
  not a cross-section of individuals with \emph{prevalent} infections.}
A final population worth distinguishing is
the "sampled population" in a given cross-sectional survey,
which we assume intends to be representative of the active population,
but might undersample some individuals,
\eg those in the early years of sex work,
due to delays in self-identification as a sex worker
\cite{Mishra2012,Cheuk2020}.
\begin{figure}
  \centering\includegraphics[scale=1]{diag.pop}
  \caption{Diagram of periods selling/buying sex in an open population at equilibrium}
  \label{fig:diag.pop}
  \floatfoot{Colours ---
    dark green: past sex work among active population (observed);
    light green: future sex work among active population (unobserved);
    gray: past/future total sex work periods among inactive population (unobserved).}
\end{figure}
